\documentclass[11pt]{article}

\usepackage[a4paper,margin=1in]{geometry}
\usepackage{amsmath,amssymb,amsthm,mathtools}
\usepackage{mathrsfs}
\usepackage[T1]{fontenc}
\usepackage{lmodern}
\usepackage{microtype}

\newtheorem{lemma}{Lemma}[section]
\newtheorem{proposition}[lemma]{Proposition}
\newtheorem{remark}[lemma]{Remark}
\newtheorem{definition}[lemma]{Definition}

\newcommand{\C}{\mathbb C}
\newcommand{\R}{\mathbb R}

\title{Block Jacobi and Riccati Reductions for the Clean OBC Hatano--Nelson Family}
\author{}
\date{}

\begin{document}
\maketitle

\section{Setup}

Fix real numbers \(0<a<b\), and let
\[
A_n=\operatorname{tridiag}(a,0,b)\in\R^{n\times n}.
\]
For \(z=x+iy\in\C\), define
\[
M_n(x,y)\coloneqq (zI-A_n)^*(zI-A_n).
\]
For a real parameter \(s\), introduce the shifted matrix
\[
H_n(s;x,y)\coloneqq M_n(x,y)-(x^2+y^2+s)I_n.
\]
Thus
\[
\det H_n(s;x,y)=0
\quad\Longleftrightarrow\quad
x^2+y^2+s
\ \text{is an eigenvalue of }M_n(x,y).
\]

We also fix the abbreviations
\[
\eta\coloneqq ab,
\qquad
c\coloneqq (a+b)x-i(b-a)y,
\qquad
\xi\coloneqq \Re c=(a+b)x,
\qquad
\zeta\coloneqq \Im c=-(b-a)y.
\]
Then
\[
-c=-a(x+iy)-b(x-iy),
\qquad
-\bar c=-a(x-iy)-b(x+iy).
\]

The \(2\times 2\) matrices that will appear repeatedly are
\[
A\coloneqq
\begin{pmatrix}
\eta & 0\\
-c & \eta
\end{pmatrix},
\]
\[
B_L\coloneqq
\begin{pmatrix}
a^2-s & -c\\
-\bar c & a^2+b^2-s
\end{pmatrix},
\]
\[
B\coloneqq
\begin{pmatrix}
a^2+b^2-s & -c\\
-\bar c & a^2+b^2-s
\end{pmatrix},
\qquad
B_R\coloneqq
\begin{pmatrix}
a^2+b^2-s & -c\\
-\bar c & b^2-s
\end{pmatrix}.
\]
For the odd case we also use
\[
d\coloneqq
\binom{\eta}{-c},
\qquad
e\coloneqq b^2-s.
\]

\section{Block Jacobi forms}

\begin{lemma}[Even case: \(n=2m\)]\label{lem:even-block}
Let \(m\ge 2\), and set
\[
w_k\coloneqq \binom{u_{2k-1}}{u_{2k}},
\qquad
1\le k\le m,
\]
for a vector \(u=(u_1,\dots,u_{2m})^T\in\C^{2m}\). Then
\[
H_{2m}(s;x,y)u=0
\]
is equivalent to the \(2\times 2\) block Jacobi system
\[
B_L w_1 + A w_2 = 0,
\]
\[
A^* w_{k-1} + B w_k + A w_{k+1} = 0,
\qquad
2\le k\le m-1,
\]
\[
A^* w_{m-1} + B_R w_m = 0.
\]
\end{lemma}

\begin{proof}
This is a direct regrouping of the rows of \(H_{2m}(s;x,y)\) into pairs
\((2k-1,2k)\), using the explicit form of \(H_{2m}(s;x,y)\):
\[
H_{2m}(s;x,y)=
\begin{pmatrix}
a^2-s & -c & \eta \\
-\bar c & a^2+b^2-s & -c & \eta \\
\eta & -\bar c & a^2+b^2-s & -c & \eta \\
& \eta & -\bar c & \ddots & \ddots & \ddots \\
&& \ddots & \ddots & a^2+b^2-s & -c & \eta \\
&&& \eta & -\bar c & a^2+b^2-s & -c \\
&&&& \eta & -\bar c & b^2-s
\end{pmatrix}.
\]
The first block row gives \(B_L w_1+A w_2=0\), the interior block rows give
\(A^*w_{k-1}+Bw_k+Aw_{k+1}=0\), and the last block row gives
\(A^*w_{m-1}+B_Rw_m=0\).
\end{proof}

\begin{lemma}[Odd case: \(n=2m+1\)]\label{lem:odd-block}
Let \(m\ge 2\), and set
\[
w_k\coloneqq \binom{u_{2k-1}}{u_{2k}},
\qquad
1\le k\le m,
\qquad
t\coloneqq u_{2m+1},
\]
for a vector \(u=(u_1,\dots,u_{2m+1})^T\in\C^{2m+1}\). Then
\[
H_{2m+1}(s;x,y)u=0
\]
is equivalent to the mixed block/scalar system
\[
B_L w_1 + A w_2 = 0,
\]
\[
A^* w_{k-1} + B w_k + A w_{k+1} = 0,
\qquad
2\le k\le m-1,
\]
\[
A^* w_{m-1} + B w_m + d\,t = 0,
\]
\[
d^* w_m + e\,t = 0.
\]
\end{lemma}

\begin{proof}
Again this is obtained by grouping the rows of \(H_{2m+1}(s;x,y)\) into pairs
\((2k-1,2k)\), together with the final scalar row \(2m+1\). The only difference
from the even case is that the odd chain has one extra site at the right endpoint,
which appears through the terminal coupling vector \(d=(\eta,-c)^T\) and scalar
terminal block \(e=b^2-s\).
\end{proof}

\section{Tail Schur complements and determinant factorizations}

\begin{proposition}[Even tail Schur complements]\label{prop:even-schur}
Assume \(m\ge 2\). On the Zariski-open set where all inverses below exist, define
recursively
\[
\Sigma_m^{(e)}\coloneqq B_R,
\qquad
\Sigma_k^{(e)}\coloneqq B_k-A\bigl(\Sigma_{k+1}^{(e)}\bigr)^{-1}A^*,
\qquad
k=m-1,\dots,1,
\]
where
\[
B_1\coloneqq B_L,
\qquad
B_k\coloneqq B
\quad (2\le k\le m-1).
\]
Then
\[
\det H_{2m}(s;x,y)=\prod_{k=1}^m \det \Sigma_k^{(e)}.
\]
In particular,
\[
\det H_{2m}(s;x,y)=0
\quad\Longleftrightarrow\quad
\det \Sigma_1^{(e)}(s;x,y)=0
\]
whenever \(\det \Sigma_k^{(e)}\neq 0\) for \(2\le k\le m\). The product formula
extends to all \(s\in\R\) by polynomial continuation.
\end{proposition}

\begin{proof}
Apply successive Schur complements starting from the rightmost block of the
block tridiagonal system in Lemma~\ref{lem:even-block}. Each step replaces the
tail block \(\Sigma_{k+1}^{(e)}\) by the effective block
\[
\Sigma_k^{(e)}=B_k-A(\Sigma_{k+1}^{(e)})^{-1}A^*.
\]
At each step, the determinant picks up a factor \(\det \Sigma_{k+1}^{(e)}\).
This yields
\[
\det H_{2m}(s;x,y)=\prod_{k=1}^m \det \Sigma_k^{(e)}
\]
on the open set where the recursion is defined. Since both sides are polynomial
functions of \(s\), the identity extends to all \(s\).
\end{proof}

\begin{proposition}[Odd tail Schur complements]\label{prop:odd-schur}
Assume \(m\ge 2\). On the open set \(s\neq b^2\), define
\[
\Sigma_m^{(o)}\coloneqq B-d\,e^{-1}d^*,
\]
and then recursively
\[
\Sigma_k^{(o)}\coloneqq B_k-A\bigl(\Sigma_{k+1}^{(o)}\bigr)^{-1}A^*,
\qquad
k=m-1,\dots,1,
\]
with
\[
B_1\coloneqq B_L,
\qquad
B_k\coloneqq B
\quad (2\le k\le m-1).
\]
Then
\[
\det H_{2m+1}(s;x,y)
=
(b^2-s)\prod_{k=1}^m \det \Sigma_k^{(o)}.
\]
In particular,
\[
\det H_{2m+1}(s;x,y)=0
\quad\Longleftrightarrow\quad
\det \Sigma_1^{(o)}(s;x,y)=0
\]
whenever \(b^2-s\neq 0\) and \(\det \Sigma_k^{(o)}\neq 0\) for \(2\le k\le m\).
The product formula extends to all \(s\in\R\) by polynomial continuation.
\end{proposition}

\begin{proof}
By Lemma~\ref{lem:odd-block}, the odd case has a scalar terminal block \(e=b^2-s\).
On the set \(e\neq 0\), eliminate the scalar tail \(t\) from the last two equations:
\[
t=-e^{-1}d^*w_m.
\]
Substituting this into the last block equation yields the effective terminal block
\[
\Sigma_m^{(o)}=B-d\,e^{-1}d^*.
\]
The remaining \(2\times 2\) block chain is then identical in form to the even case,
so successive Schur complements give
\[
\det H_{2m+1}(s;x,y)
=
e\prod_{k=1}^m \det \Sigma_k^{(o)}.
\]
As before, both sides are polynomial in \(s\), so the identity extends to all \(s\).
\end{proof}

\section{Scalar Riccati recursions}

\begin{proposition}[Even scalar Riccati recursion]\label{prop:even-scalar}
Write
\[
\Sigma_k^{(e)}=
\begin{pmatrix}
u_k^{(e)} & \beta_k^{(e)}\\
\overline{\beta_k^{(e)}} & v_k^{(e)}
\end{pmatrix},
\qquad
\Delta_k^{(e)}\coloneqq
u_k^{(e)}v_k^{(e)}-\bigl|\beta_k^{(e)}\bigr|^2.
\]
Then the terminal data are
\[
u_m^{(e)}=a^2+b^2-s,
\qquad
\beta_m^{(e)}=-c,
\qquad
v_m^{(e)}=b^2-s,
\]
and for \(k=m-1,\dots,1\) one has
\[
u_k^{(e)}
=
\alpha_k-\frac{\eta^2 v_{k+1}^{(e)}}{\Delta_{k+1}^{(e)}},
\]
\[
\beta_k^{(e)}
=
-c+\frac{\eta\bigl(\eta\,\beta_{k+1}^{(e)}+v_{k+1}^{(e)}\bar c\bigr)}
{\Delta_{k+1}^{(e)}},
\]
\[
v_k^{(e)}
=
\delta_k-\frac{
\eta^2 u_{k+1}^{(e)}
+\eta\bigl(\beta_{k+1}^{(e)}c+\overline{\beta_{k+1}^{(e)}}\,\bar c\bigr)
+|c|^2 v_{k+1}^{(e)}
}{\Delta_{k+1}^{(e)}},
\]
where
\[
(\alpha_1,\delta_1)=(a^2-s,\ a^2+b^2-s),
\qquad
(\alpha_k,\delta_k)=(a^2+b^2-s,\ a^2+b^2-s)
\quad (2\le k\le m-1).
\]
\end{proposition}

\begin{proof}
For a Hermitian matrix
\[
\Sigma=
\begin{pmatrix}
u & \beta\\
\bar\beta & v
\end{pmatrix},
\qquad
\Delta=uv-|\beta|^2,
\]
one has
\[
\Sigma^{-1}
=
\frac1\Delta
\begin{pmatrix}
v & -\beta\\
-\bar\beta & u
\end{pmatrix}.
\]
Now insert this into
\[
\Sigma_k^{(e)}=B_k-A\bigl(\Sigma_{k+1}^{(e)}\bigr)^{-1}A^*,
\]
with
\[
A=
\begin{pmatrix}
\eta & 0\\
-c & \eta
\end{pmatrix}.
\]
A direct multiplication yields
\[
A
\begin{pmatrix}
v & -\beta\\
-\bar\beta & u
\end{pmatrix}
A^*
=
\begin{pmatrix}
\eta^2 v &
-\eta(\eta\beta+v\bar c)\\[0.4ex]
-\eta(\eta\bar\beta+cv) &
\eta^2u+\eta(\beta c+\bar\beta\,\bar c)+|c|^2v
\end{pmatrix}.
\]
Subtracting this from \(B_k\) gives the stated formulas.
\end{proof}

\begin{proposition}[Odd scalar Riccati recursion]\label{prop:odd-scalar}
Write
\[
\Sigma_k^{(o)}=
\begin{pmatrix}
u_k^{(o)} & \beta_k^{(o)}\\
\overline{\beta_k^{(o)}} & v_k^{(o)}
\end{pmatrix},
\qquad
\Delta_k^{(o)}\coloneqq
u_k^{(o)}v_k^{(o)}-\bigl|\beta_k^{(o)}\bigr|^2.
\]
On the set \(s\neq b^2\), the terminal data are
\[
u_m^{(o)}
=
a^2+b^2-s-\frac{\eta^2}{b^2-s},
\]
\[
\beta_m^{(o)}
=
-c+\frac{\eta\,\bar c}{b^2-s},
\]
\[
v_m^{(o)}
=
a^2+b^2-s-\frac{|c|^2}{b^2-s},
\]
and for \(k=m-1,\dots,1\) one has exactly the same Riccati recursion as in the even case:
\[
u_k^{(o)}
=
\alpha_k-\frac{\eta^2 v_{k+1}^{(o)}}{\Delta_{k+1}^{(o)}},
\]
\[
\beta_k^{(o)}
=
-c+\frac{\eta\bigl(\eta\,\beta_{k+1}^{(o)}+v_{k+1}^{(o)}\bar c\bigr)}
{\Delta_{k+1}^{(o)}},
\]
\[
v_k^{(o)}
=
\delta_k-\frac{
\eta^2 u_{k+1}^{(o)}
+\eta\bigl(\beta_{k+1}^{(o)}c+\overline{\beta_{k+1}^{(o)}}\,\bar c\bigr)
+|c|^2 v_{k+1}^{(o)}
}{\Delta_{k+1}^{(o)}},
\]
with
\[
(\alpha_1,\delta_1)=(a^2-s,\ a^2+b^2-s),
\qquad
(\alpha_k,\delta_k)=(a^2+b^2-s,\ a^2+b^2-s)
\quad (2\le k\le m-1).
\]
\end{proposition}

\begin{proof}
The terminal data come from the explicit matrix
\[
\Sigma_m^{(o)}=B-d\,e^{-1}d^*
=
\begin{pmatrix}
a^2+b^2-s-\dfrac{\eta^2}{b^2-s}
&
-c+\dfrac{\eta\,\bar c}{b^2-s}\\[2ex]
-\bar c+\dfrac{\eta\,c}{b^2-s}
&
a^2+b^2-s-\dfrac{|c|^2}{b^2-s}
\end{pmatrix}.
\]
Once this terminal block is fixed, the proof is identical to that of
Proposition~\ref{prop:even-scalar}.
\end{proof}

\section{Lorentz-cone form}

\begin{definition}[Lorentz parametrization of \(2\times 2\) Hermitian matrices]
For \(X=(m,\chi,\psi,z)^T\in\R^4\), define
\[
\mathscr S(X)\coloneqq
\begin{pmatrix}
m+z & \chi-i\psi\\
\chi+i\psi & m-z
\end{pmatrix}.
\]
Also define the quadratic form
\[
Q(X)\coloneqq m^2-\chi^2-\psi^2-z^2,
\]
and the involution
\[
J\coloneqq \operatorname{diag}(1,-1,-1,-1).
\]
Then
\[
\det \mathscr S(X)=Q(X),
\qquad
\mathscr S(X)^{-1}=\frac{1}{Q(X)}\,\mathscr S(JX)
\]
whenever \(Q(X)\neq 0\).
\end{definition}

\begin{proposition}[Common Lorentz-cone Riccati map]\label{prop:lorentz}
Define the real \(4\times 4\) matrix
\[
\mathcal L_c\coloneqq
\begin{pmatrix}
\eta^2+\dfrac{|c|^2}{2} & -\eta \xi & -\eta \zeta & \dfrac{|c|^2}{2}\\[1.2ex]
-\eta \xi & \eta^2 & 0 & -\eta \xi\\[1.2ex]
-\eta \zeta & 0 & \eta^2 & -\eta \zeta\\[1.2ex]
-\dfrac{|c|^2}{2} & \eta \xi & \eta \zeta & \eta^2-\dfrac{|c|^2}{2}
\end{pmatrix}.
\]
Then
\[
A\,\mathscr S(X)\,A^*=\mathscr S(\mathcal L_c X)
\qquad\text{for every }X\in\R^4.
\]

Now let \(\Sigma_k^{(\#)}=\mathscr S(X_k^{(\#)})\), where \(\#\in\{e,o\}\) denotes the even or odd case.
Then the Riccati recursions of Propositions~\ref{prop:even-scalar} and \ref{prop:odd-scalar}
become
\[
X_k^{(\#)}
=
X_{B_k}-\frac{1}{Q(X_{k+1}^{(\#)})}\,\mathcal L_c\,JX_{k+1}^{(\#)},
\qquad
k=m-1,\dots,1,
\]
where
\[
X_{B_L}=
\left(
a^2+\frac{b^2}{2}-s,\ -\xi,\ \zeta,\ -\frac{b^2}{2}
\right)^T,
\]
\[
X_B=
\left(
a^2+b^2-s,\ -\xi,\ \zeta,\ 0
\right)^T.
\]
The terminal data are
\[
\Sigma_m^{(e)}=B_R
\quad\Longleftrightarrow\quad
X_m^{(e)}=
\left(
\frac{a^2+2b^2}{2}-s,\ -\xi,\ \zeta,\ \frac{a^2}{2}
\right)^T,
\]
and, on \(s\neq b^2\),
\[
\Sigma_m^{(o)}=B-d\,e^{-1}d^*
\quad\Longleftrightarrow\quad
X_m^{(o)}=
\begin{pmatrix}
a^2+b^2-s-\dfrac{\eta^2+|c|^2}{2(b^2-s)}\\[0.6ex]
-\xi+\dfrac{\eta\xi}{b^2-s}\\[0.6ex]
\zeta+\dfrac{\eta\zeta}{b^2-s}\\[0.6ex]
\dfrac{|c|^2-\eta^2}{2(b^2-s)}
\end{pmatrix}.
\]
\end{proposition}

\begin{proof}
The identity
\[
A\,\mathscr S(X)\,A^*=\mathscr S(\mathcal L_c X)
\]
is a direct computation. Since
\[
\mathscr S(X)^{-1}=\frac{1}{Q(X)}\,\mathscr S(JX),
\]
one gets
\[
A\,\mathscr S(X)^{-1}A^*
=
\frac{1}{Q(X)}\,A\,\mathscr S(JX)\,A^*
=
\frac{1}{Q(X)}\,\mathscr S(\mathcal L_c JX).
\]
Substituting this into
\[
\Sigma_k^{(\#)}=B_k-A\bigl(\Sigma_{k+1}^{(\#)}\bigr)^{-1}A^*
\]
and identifying both sides via the map \(\mathscr S\) yields the stated Lorentz-cone recursion.
The terminal vectors are obtained by decomposing the explicit terminal matrices
\(\Sigma_m^{(e)}\) and \(\Sigma_m^{(o)}\) in the basis \(\mathscr S(\cdot)\).
\end{proof}

\begin{remark}
The results above show that, after odd/even blocking, the spectral problem for
\(H_n(s;x,y)\) is governed by a rational self-map of the Lorentz cone
\[
Q(X)>0,\qquad m>0,
\]
that is, by a \(4\)-real-dimensional Riccati dynamics on \(2\times 2\) Hermitian
tail Schur complements. In the even and odd cases, the interior recursion is the
same; only the terminal data differ.
\end{remark}

\begin{remark}
For \(n=2\) and \(n=3\), the above constructions collapse to one-step problems and
recover the explicit scalar analyses of those dimensions.
\end{remark}

\begin{proposition}[Closure of the terminal cocycle contribution]
Let
\[
\widehat H_n(\Lambda;x,y)=M_n(x,y)-\Lambda I,
\qquad
q=b-a,
\qquad
\eta=ab,
\qquad
c=(a+b)x-iqy.
\]
Let \(\Sigma_k(\Lambda,y)\) be the blocked tail Schur complements, and define
\[
S_k:=-\partial_\Lambda\Sigma_k,
\qquad
P_k:=2y\,S_k-\partial_y\Sigma_k.
\]
Let \(w_1\in\ker\Sigma_1\setminus\{0\}\), and propagate the nullchain by
\[
w_{k+1}=-\Sigma_{k+1}^{-1}A^*w_k.
\]
Write
\[
j_k:=\Im(\overline{w_{k,1}}\,w_{k,2}).
\]
Assume \(y>0\). Then the terminal contribution is strictly positive.

\smallskip

\noindent
\textup{(i) Even case \(n=2m\).} One has
\[
\Sigma_m=B_R,
\qquad
S_m=I_2,
\qquad
P_m=
\begin{pmatrix}
0 & -iq\\
iq & 0
\end{pmatrix},
\]
and therefore
\[
w_m^*P_mw_m=2q\,j_m>0.
\]

\smallskip

\noindent
\textup{(ii) Odd case \(n=2m+1\).} Set
\[
d_0:=\binom{\eta}{-c},
\qquad
e:=x^2+y^2+b^2-\Lambda.
\]
Then
\[
\Sigma_m=B-d_0e^{-1}d_0^*,
\qquad
S_m=I_2+e^{-2}d_0d_0^*,
\]
and
\[
P_m=
\begin{pmatrix}
0 & -iq\left(1+\dfrac{\eta}{e}\right)\\[1.2ex]
iq\left(1+\dfrac{\eta}{e}\right) & \dfrac{2q^2y}{e}
\end{pmatrix}.
\]
Moreover \(e>0\) at \(\Lambda=\lambda_{\min}(M_n(x,y))\), hence
\[
w_m^*P_mw_m
=
2q\left(1+\frac{\eta}{e}\right)j_m
+
\frac{2q^2y}{e}|w_{m,2}|^2
>0.
\]

\smallskip

Consequently, if
\[
P_k=\Pi_k+R_kP_{k+1}R_k^*,
\qquad
R_k=A\Sigma_{k+1}^{-1},
\]
and the interior terms satisfy \(w_k^*\Pi_kw_k>0\) for \(k=1,\dots,m-1\), then
\[
w_1^*P_1w_1
=
\sum_{k=1}^{m-1} w_k^*\Pi_kw_k+w_m^*P_mw_m
>0.
\]
\end{proposition}

\begin{proof}
The even case is immediate from
\[
P_m=2yI_2-\partial_yB_R=
\begin{pmatrix}
0 & -iq\\
iq & 0
\end{pmatrix}.
\]
For the odd case, differentiate
\[
\Sigma_m=B-d_0e^{-1}d_0^*
\]
with respect to \(\Lambda\) and \(y\), using
\[
\partial_\Lambda(e^{-1})=e^{-2},
\qquad
\partial_y(e^{-1})=-\frac{2y}{e^2},
\qquad
\partial_y d_0=\binom{0}{iq},
\]
to obtain the stated formulas for \(S_m\) and \(P_m\). The positivity of the
quadratic forms then follows from \(q>0\), \(y>0\), \(e>0\), and \(j_m>0\).
\end{proof}

\end{document}
